\section{Dataset and Problem Definition}

% Progress Report: 3.1 and 3.2 in full, 3.3 to 3.6 as the PLANNED protocol (one short paragraph each
% or less). 3.7 and 3.8 are Final Report only.

\subsection{Task Definition}
% WRITE: the input is a continuous RGB stream from a single ordinary webcam; the output is the
% recognised gesture (or ``no gesture''). Three tasks: static classification, dynamic
% classification, temporal segmentation. Name the difficulties that will be evaluated
% (lighting, background, camera distance, users, intentional versus incidental movement) and the
% constraints (real time, ordinary webcam, no extra hardware).
\todo{task definition}

\subsection{Gesture Set}
% The gesture set below is DECIDED. Commands are not defined yet: the system only recognises the
% gestures for now. The label names are the ones in dataset_config.py.
\begin{table}[h]
\centering
\begin{tabular}{@{}lll p{6.2cm}@{}}
\toprule
Label & Type & Hand & Description \\
\midrule
\texttt{fist}        & static  & closed hand & held still \\
\texttt{open\_palm}  & static  & open hand & fingers extended, held still \\
\texttt{thumb\_up}   & static  & fist & thumb pointing up, held still \\
\texttt{thumb\_down} & static  & fist & thumb pointing down, held still \\
\texttt{point}       & static  & index finger & index finger extended, other fingers closed \\
\texttt{swipe\_left}  & dynamic & open palm & out-and-back movement; first stroke towards the user's left \\
\texttt{swipe\_right} & dynamic & open palm & out-and-back movement; first stroke towards the user's right \\
\texttt{no\_gesture}  & both    & any & negative class: poses and movements that are not one of the gestures \\
\bottomrule
\end{tabular}
\caption{Gesture set. \todo{decide whether to add the commands (target device) before the Final Report}}
\label{tab:gestures}
\end{table}
% WRITE: one sentence per gesture on why it was chosen (visually distinct, easy to perform, not
% confusable). Explain that the open palm is used both as a static gesture and for swipes, and
% that the two are told apart by movement (see Section on segmentation).

\subsection{Data Collection Protocol}
% PLANNED protocol (Progress Report).
% FACTS from the project, check before use:
%  - camera 640 x 480, about 20 frames per second while recording;
%  - only the 21 hand landmarks are stored, not the images; clips are labelled with user,
%    gesture, background, lighting and distance;
%  - static clips: 2 seconds, hand held completely still; swipe and no_gesture clips: 3 seconds;
%  - swipe: wait about half a second after the countdown, make the out-and-back movement once at a
%    steady speed with a large amplitude (about one hand size or more), then keep the hand still;
%  - participants give consent; raw video is not kept.
\todo{protocol}

\subsection{Participants and Recording Conditions}
% PLANNED (Progress Report). Fill in the real numbers.
\begin{table}[h]
\centering
\begin{tabular}{@{}lll@{}}
\toprule
Factor & Levels & Planned \\
\midrule
Users & \todo{how many} & \todo{} \\
Background & plain, cluttered & \todo{} \\
Lighting & bright, dim & \todo{} \\
Camera distance & near, far & \todo{} \\
\bottomrule
\end{tabular}
\caption{Recording conditions.}
\end{table}

\subsection{Dataset Statistics}
% Numbers come from validate_dataset.py (clips per user, gesture and condition).
% Progress Report: planned numbers, or the pilot numbers clearly marked as pilot.
\todo{statistics table}

\subsection{Train, Validation and Test Split}
% WRITE: the split is by PERSON, never by frame or clip, because frames of one clip are nearly
% identical and would leak into the test set. Test users are not used while choosing thresholds or
% models; they are used once, at the end. With few users, leave-one-user-out cross-validation.
\todo{split}

\iffinal
\subsection{Preprocessing and Augmentation}
% WRITE: smoothing, normalisation, stable-frame selection; any augmentation (for example random
% rotation of the training frames).
\todo{preprocessing and augmentation}

\subsection{Data Limitations}
% WRITE: known biases and gaps (number of users, one camera, one environment, ...).
\todo{limitations}
\fi
